\section{How Much of the Gain Is Real?}
\label{sec:heldout}

\S\ref{sec:learned} establishes that the policy changed in ways MI would not endorse. It does
not establish how much of the \emph{measured} improvement those changes account for, because
every number so far comes from the model that also supplied the reward. To separate the two we
re-scored the entire study with a second grader.

\subsection{A decoupled second grader}
\heldoutjudge{} re-scored all \textbf{16{,}896} evaluation cells of this comparison --- 8
instruments $\times$ 22 model states $\times$ 96 conversations, complete, no missing cells (part
of a larger sweep covering all 39 model states in the project, $29{,}952$ cells). It is decoupled from
training in three ways that matter: a different model family, it never played the patient, and
it never produced a reward. Rubrics were held structurally identical, verified by an automated
parity check before any scoring (Appendix~\ref{app:measurement}).

\textbf{This is a held-out grader, not a second rater.} The primary oracle \emph{was} the
training signal, so the relation is train-versus-test, not inter-annotator agreement. We
therefore never average the two --- the level offset is $1.2$--$1.7$ points and
model-dependent, so averaging would silently shrink every effect by an arm-dependent amount.
Only contrasts are compared.

\subsection{The rankings survive; the levels do not}
Across all \textbf{1{,}848} pairwise arm\,$\times$\,instrument contrasts, $88.3\%$ keep their
sign under the held-out judge, rising to $94.1\%$ at $|\Delta|\geq0.10$ and $98.9\%$ at
$\geq0.50$: the judges disagree only about differences too small to claim. Just $1.2$--$6.9\%$
of arm-mean variance sits in the arm\,$\times$\,judge interaction --- the one component that
could invalidate a ranking (Appendix~\ref{app:measurement}). \heldoutjudge{} is systematically
harsher, which cancels in every contrast.

\textbf{MITI is the exception}: $3.6\%$ between-arm signal, dependability $0.65$, sign preserved
on only $77.5\%$ of contrasts, and the sole instrument still disagreeing at claimable effect
sizes. We therefore treat every MITI-based statement in \S\ref{sec:gains} --- including the
competency thresholds --- as provisional. Q1, Q2, PCT and MICI are unaffected.

\subsection{Gain retention}
What matters is not whether the graders rank arms alike but how much of each arm's improvement
the held-out grader will credit. We define \textbf{gain retention} for a model state $m$ on
instrument $i$ as
\[
R_i(m)\;=\;\frac{\Delta^{\text{held-out}}_i(m)}{\Delta^{\text{trained-against}}_i(m)},
\]
with $\Delta_i(m)$ the persona-paired improvement of $m$ over base. $R=1$ credits the full
gain, $R=0$ none of it. Since both terms are differences from the same reference, the level
offset cancels; intervals are bootstrapped over the 96 personas.

\begin{table}[t]
\centering\small
\setlength{\tabcolsep}{3pt}
\begin{tabular}{@{}llrrl@{}}
\toprule
& & $\Delta$\,train & $\Delta$\,held & \textbf{retention} \\
& & \multicolumn{1}{c}{\footnotesize gpt-4o-mini} & \multicolumn{1}{c}{\footnotesize Haiku} & \footnotesize [95\% CI] \\
\midrule
\multirow{4}{*}{\textbf{Q1}}
 & PTO\,@10  & 1.22 & 0.97 & \textbf{0.80} [0.68, 0.93] \\
 & GRPO\,@10 & 0.68 & 0.19 & \textbf{0.28} [0.06, 0.43] \\
 & PTO\,@8   & 1.18 & 1.05 & 0.89 [0.76, 1.04] \\
 & GRPO\,@8  & 1.01 & 0.65 & 0.64 [0.52, 0.78] \\
\midrule
\multirow{4}{*}{\textbf{Q2}}
 & PTO\,@10  & 1.30 & 1.11 & 0.85 [0.74, 0.98] \\
 & GRPO\,@10 & 0.82 & 0.66 & 0.81 [0.68, 0.96] \\
 & PTO\,@8   & 1.26 & 1.08 & 0.86 [0.74, 0.99] \\
 & GRPO\,@8  & 1.15 & 0.92 & 0.80 [0.68, 0.95] \\
\bottomrule
\end{tabular}
\caption{Gain retention $=\Delta_\text{held-out}/\Delta_\text{trained-against}$, where
$\Delta$ is the persona-paired improvement over the base policy ($n=96$). \textbf{Q1 separates
the arms with non-overlapping intervals at the matched iteration-10 endpoint}; at GRPO's
iteration-8 peak PTO still leads but the intervals overlap marginally, so the peak-vs-peak
claim is an ordering rather than a disjoint contrast (\S\ref{sec:heldout}). \textbf{Q2 does not
separate at all} --- all four intervals sit in a $0.80$--$0.86$ band, the control showing that
retention is not merely a stingier grader.}
\label{tab:retention}
\end{table}


\begin{figure}[t]
\centering
\includegraphics[width=\columnwidth]{figures/multijudge_retention_trajectory.png}
\caption{Gain retention by iteration, every instrument where the ratio is estimable (MICI's
base $\Delta$ is near zero, so only iterations 8--10 survive). Blue PTO, orange GRPO;
$1.0$ = the held-out judge credits the full gain. \textbf{The arms separate on Q1 and MITI and
track each other on the alliance and satisfaction scales.} In the Q1 panel they are
indistinguishable for three iterations, after which PTO holds $0.80$--$0.98$ while GRPO decays
in trend to $0.28$ --- visible well before the outcome curve (Figure~\ref{fig:trajectory})
turns over at iteration~8. Retention above $1.0$ on WAI-SR and CSQ-8 means the held-out judge
credits \emph{more} there than the primary, which is why the diagnostic is the between-arm
comparison within an instrument and not the absolute level.}
\label{fig:retention}
\end{figure}

At the matched endpoint (Table~\ref{tab:retention}), \textbf{PTO retains $0.80$
[$0.68$, $0.93$] of its Q1 gain and GRPO $0.28$ [$0.06$, $0.43$]} --- non-overlapping
intervals. Stated in the units of the instrument: under a grader that never graded during
training, GRPO's net ten-iteration Q1 improvement is ${\approx}0.19$ points, not the
${\approx}0.68$ its own reward reports.

\paragraph{Does it survive a best-vs-best comparison?} Only partly, and we state the limit
precisely. The intervals are non-overlapping at iterations 9--10 \emph{only}. At GRPO's own
iteration-8 peak PTO still leads ($0.89$ vs $0.64$), but the intervals overlap marginally
($[0.76,1.04]$ vs $[0.52,0.78]$), as do PTO@10 against GRPO@8 ($[0.68,0.93]$ vs $[0.52,0.78]$).
\textbf{What is robust is the ordering, not any single disjoint pair}: PTO's Q1 retention
exceeds GRPO's at \emph{every} iteration from 4 onward (Figure~\ref{fig:retention}), with
monotonically diverging trends. A reader who insists on peak-vs-peak should take the retention
result as a consistent ordering whose endpoint separates, not as one significant contrast.

\paragraph{Q2 is a control, and the separation is instrument-specific.} Every Q2 retention
interval overlaps, in a narrow $0.80$--$0.86$ band across both arms and endpoints --- mild scale
compression affecting both equally. That matters because it shows retention is not simply
measuring ``the second grader is stingier''; if it were, Q1 and Q2 would behave alike. They do
not. Figure~\ref{fig:retention} shows the separation is confined to Q1 and MITI, while WAI-SR,
CSQ-8 and MI-SAT track each other in both arms (and sit above $1.0$, the held-out judge crediting
\emph{more} than the primary). We read this as diagnostic rather than dilutive: Q1 asks
\emph{was this helpful, did I learn, will I act on it} --- where warmth substitutes for
substance most easily --- whereas the alliance and satisfaction scales reward warmth directly,
and so have less to withdraw.

\paragraph{Onset, not endpoint.} Retention is a curve, not a final number. Both arms sit at
$0.84$--$0.97$ through iteration~3; GRPO then declines steadily before collapsing at iterations
9--10 ($0.03$, $0.28$), while PTO stays in $0.80$--$0.98$ throughout. Two consequences.
Retention separates the arms from iteration~4, whereas the reward curve does not distinguish
them until GRPO turns over at iteration~8 --- a monitor computed on data one already has would
have flagged this four iterations earlier. And GRPO's outcome regression and its retention
collapse are the same phenomenon seen twice.

\paragraph{Why retention rather than MI-inconsistency.} Retention needs no cross-judge agreement
on level or rate --- only that each grader is internally consistent across the states it
compares --- whereas MICI's agreement is weak ($r=0.20$--$0.55$). Hence retention is
load-bearing here and MICI corroborating.
